\documentclass[10pt]{article}
\usepackage[T1]{fontenc}
\usepackage[utf8]{inputenc}
\usepackage[french]{babel}
\usepackage{amsmath,amssymb,mathrsfs,mathpazo}
\usepackage{geometry}
\geometry{a4paper,margin=2cm}
\usepackage{graphicx,array,tabularx,tabularray,colortbl,multirow,multicol,booktabs,enumitem,pifont,ulem,fancybox,diagbox,cancel}
\usepackage{pstricks,pstricks-add,pst-plot,pst-tree,pst-node,pst-text,pst-func,pst-math,pst-bspline,pst-3dplot,pst-eucl,pst-solides3d}
\let\pscalculate\undefined
\usepackage{pst-bezier}
\usepackage{tikz,pgfplots,tkz-tab}
\usetikzlibrary{arrows,arrows.meta,patterns,calc,positioning,trees,decorations.markings,shapes.geometric,shapes.misc,shapes.arrows,intersections,decorations.pathmorphing,decorations.pathreplacing,backgrounds,shadows,babel}
\pgfplotsset{compat=1.18}
\usepgfplotslibrary{fillbetween}
\usepackage[np]{numprint}
\usepackage{listings,makecell,pas-tableur,twoopt}
\let\barre\undefined
\usepackage[pstricks]{tablvar}
\definecolor{darkgreen}{RGB}{0,100,0}
\providecommand{\red}{\color{red}}
\providecommand{\blue}{\color{blue}}
\providecommand{\green}{\color{green}}
\providecommand{\black}{\color{black}}
\providecommand{\white}{\color{white}}
\providecommand{\gray}{\color{gray}}
\providecommand{\cyan}{\color{cyan}}
\providecommand{\magenta}{\color{magenta}}
\providecommand{\yellow}{\color{yellow}}
\providecommand{\violet}{\color{violet}}
\providecommand{\orange}{\color{orange}}
\providecommand{\euro}{\texteuro}
\providecommand{\EUR}{\texteuro}
\providecommand{\ieme}{\textsuperscript{e}}
\providecommand{\degres}{\ensuremath{^\circ}}
\providecommand{\degree}{\ensuremath{^\circ}}
\providecommand{\vv}[1]{\overrightarrow{#1}}
\providecommand{\Coord}[1]{#1}
\providecommand{\decosix}{}
\providecommand{\starredbullet}{\ensuremath{\bullet}}
\pagestyle{empty}
\setlength{\parindent}{0pt}
\usepackage{adjustbox}\DeclareUnicodeCharacter{2236}{:}\DeclareUnicodeCharacter{2190}{\ensuremath{\leftarrow}}
\providecommand{\R}{\mathbb{R}}
\providecommand{\N}{\mathbb{N}}
\providecommand{\Z}{\mathbb{Z}}
\providecommand{\Q}{\mathbb{Q}}
\providecommand{\D}{\mathbb{D}}
\providecommand{\C}{\mathbb{C}}
\providecommand{\e}{{\mathrm{e}}}
\providecommand{\vect}[1]{\overrightarrow{#1}}
\providecommand{\vectt}[1]{\overrightarrow{\mathrm{#1}}}
\providecommand{\vvt}[1]{\overrightarrow{\mathrm{#1}}}
\providecommand{\np}[2][]{{#2~#1}}
\providecommand{\vv}[1]{{\overrightarrow{#1}}}
\renewcommand{\d}{{\mathrm{d}}}
\providecommand{\degres}{{^\circ}}
\providecommand{\degre}{{^\circ}}
\providecommand{\blue}{\color{blue}}
\providecommand{\red}{\color{red}}
\providecommand{\enskip}{\quad}
\providecommand{\llbracket}{\lbrack\!\lbrack}
\providecommand{\rrbracket}{\rbrack\!\rbrack}
\providecommand{\euro}{{\text{€}}}
\providecommand{\no}{{\text{n°}}}
\providecommand{\interval}[2]{{[#1;#2]}}
\providecommand{\intervFO}[2]{{[#1;#2[}}
\providecommand{\dotfill}{\ldots}

\providecommand{\ds}{\displaystyle}
\providecommand{\textding}[1]{#1}
\providecommand{\Oij}{\(\mathrm{O};\vec{\imath},\vec{\jmath}\)}
\providecommand{\Oijk}{\(\mathrm{O};\vec{\imath},\vec{\jmath},\vec{k}\)}
\providecommand{\Ouv}{\(\mathrm{O};\vec{u},\vec{v}\)}
\providecommand{\intervFF}[2]{\ensuremath{[#1;#2]}}
\renewcommand{\euro}{\ifmmode\text{\texteuro}\else\texteuro\fi}
\DeclareUnicodeCharacter{00B0}{\ifmmode^{\circ}\else\textdegree\fi}
\DeclareUnicodeCharacter{20AC}{\ifmmode\text{\texteuro}\else\texteuro\fi}
% Énoncés : Éducation nationale. Transcription et corrigés : APMEP, Denis Vergès et contributeurs. Découpage et classement : CoopMaths, Sébastien Lozano, Éric Elter et contributeurs. Adaptation de la présentation : Maths Explorer.
% CC BY-SA : https://coopmaths.fr/www/about/
% https://coopmaths.fr/alea/static/bac/2024/tex/bac_2024_06_sujet1_polynesie_3.tex
\begin{document}
\noindent\textit{APMEP / CoopMaths. CC BY-SA. Présentation adaptée par Maths Explorer.}\par\medskip
\medskip

\emph{Les probabilités demandées seront exprimées sous forme de fractions irréductibles}

\medskip

\textbf{Partie A}

\medskip

On lance trois fois de suite une pièce de monnaie bien équilibrée. On note $X$ la variable aléatoire qui compte le nombre de fois, sur les trois lancers, où la pièce est retombée du côté \og Face \fg.

\medskip

\begin{enumerate}
\item Préciser la nature et les paramètres de la loi de probabilité suivie par $X$.



\item Recopier et compléter le tableau suivant donnant la loi de probabilité de $X$
\begin{center}
\begin{tabularx}{\linewidth}{|*{5}{>{\centering \arraybackslash}X|}}\hline

$k$& 0 & 1 & 2& 3 \\ \hline
$P(X = k)$&&&&\\ \hline
\end{tabularx}
\end{center}
\end{enumerate}

\medskip

\textbf{Partie B}

\medskip

Voici les règles d'un jeu où le but est d'obtenir trois pièces du côté \og Face \fg{} en un ou deux essais :

\begin{itemize}[label=$\bullet~$]
\item On lance trois pièces équilibrées :
	\begin{itemize}[label=$\circ~$]
	\item Si les trois pièces sont tombées du côté \og Face \fg, la partie est gagnée ;
	\item Sinon, les pièces tombées du côté \og Face\fg{} sont conservées et on relance
celles tombées du côté \og Pile \fg.
	\end{itemize}
\item La partie est gagnée si on obtient trois pièces du côté \og Face \fg, sinon elle est
perdue.
\end{itemize}

On considère les évènements suivants :
\begin{itemize}[label=$\bullet~$]
\item $G$ : \og la partie est gagnée \fg.

Et pour tout entier $k$ compris entre 0 et 3, les évènements:
\item $A_k$ : \og $k$ pièces sont tombées du côté \og Face\fg{} au premier lancer \fg.
\end{itemize}

\medskip

\begin{enumerate}
\item Démontrer que $P_{A_1}(G) = \dfrac14$.
\item Recopier et compléter l'arbre pondéré ci-dessous:

\begin{center}
\pstree[treemode=R,levelsep=3cm]{\TR{}}
{\pstree{\TR{$A_0~~$}\taput{$\frac18$}}
	{\TR{$G$}
	\TR{$\overline{G}$}
	}
	\pstree{\TR{$A_1~~$}\taput{$\frac38$}}
	{\TR{$G$}
	\TR{$\overline{G}$}
	}
	\pstree{\TR{$A_2~~$}\taput{$\frac38$}}
	{\TR{$G$}
	\TR{$\overline{G}$}
	}
	\pstree{\TR{$A_3~~$}\tbput{$\frac18$}}
	{\TR{$G$}	
	}
}
\end{center}

\item Démontrer que la probabilité $p$ de gagner à ce jeu est $p = \dfrac{27}{64}$
\item La partie a été gagnée. Quelle est la probabilité qu'exactement une pièce soit tombée du côté \og Face \fg{} à la première tentative ?
\item Combien de fois faut-il jouer à ce jeu pour que la probabilité de gagner au moins une partie dépasse $0,95$ ?
\end{enumerate}

\bigskip

\end{document}
